\documentclass[10pt,a4paper]{amsart}
\usepackage{amsmath,amssymb,amsthm,graphicx}
% Retained minimal preamble for the Stacks Project chapter sites-cohomology.tex (Cohomology on Sites).
% Verbatim excerpt of the chapter preamble: Xy-pic setup (native preamble lines 14-18) and the native
% theorem declarations with the mathematical macros (native preamble lines 154-217). Copied mechanically;
% nothing added, removed or redefined. The wrapper class and the amsmath/amssymb/amsthm/graphicx base
% packages are supplied by the assembler, so the class/externaldocument part of the native preamble is omitted.
% For commutative diagrams we use Xy-pic
\usepackage[all]{xy}

% We use 2cell for 2-commutative diagrams.
\xyoption{2cell}
% Theorem environments.
%
\theoremstyle{plain}
\newtheorem{theorem}[subsection]{Theorem}
\newtheorem{proposition}[subsection]{Proposition}
\newtheorem{lemma}[subsection]{Lemma}

\theoremstyle{definition}
\newtheorem{definition}[subsection]{Definition}
\newtheorem{example}[subsection]{Example}
\newtheorem{exercise}[subsection]{Exercise}
\newtheorem{situation}[subsection]{Situation}

\theoremstyle{remark}
\newtheorem{remark}[subsection]{Remark}
\newtheorem{remarks}[subsection]{Remarks}

\numberwithin{equation}{subsection}

% Macros
%
\def\lim{\mathop{\mathrm{lim}}\nolimits}
\def\colim{\mathop{\mathrm{colim}}\nolimits}
\def\Spec{\mathop{\mathrm{Spec}}}
\def\Hom{\mathop{\mathrm{Hom}}\nolimits}
\def\Ext{\mathop{\mathrm{Ext}}\nolimits}
\def\SheafHom{\mathop{\mathcal{H}\!\mathit{om}}\nolimits}
\def\SheafExt{\mathop{\mathcal{E}\!\mathit{xt}}\nolimits}
\def\Sch{\mathit{Sch}}
\def\Mor{\mathop{\mathrm{Mor}}\nolimits}
\def\Ob{\mathop{\mathrm{Ob}}\nolimits}
\def\Sh{\mathop{\mathit{Sh}}\nolimits}
\def\NL{\mathop{N\!L}\nolimits}
\def\CH{\mathop{\mathrm{CH}}\nolimits}
\def\proetale{{pro\text{-}\acute{e}tale}}
\def\etale{{\acute{e}tale}}
\def\QCoh{\mathit{QCoh}}
\def\Ker{\mathop{\mathrm{Ker}}}
\def\Im{\mathop{\mathrm{Im}}}
\def\Coker{\mathop{\mathrm{Coker}}}
\def\Coim{\mathop{\mathrm{Coim}}}

% Boxtimes
%
\DeclareMathSymbol{\boxtimes}{\mathbin}{AMSa}{"02}

%
% Macros for moduli stacks/spaces
%
\def\QCohstack{\mathcal{QC}\!\mathit{oh}}
\def\Cohstack{\mathcal{C}\!\mathit{oh}}
\def\Spacesstack{\mathcal{S}\!\mathit{paces}}
\def\Quotfunctor{\mathrm{Quot}}
\def\Hilbfunctor{\mathrm{Hilb}}
\def\Curvesstack{\mathcal{C}\!\mathit{urves}}
\def\Polarizedstack{\mathcal{P}\!\mathit{olarized}}
\def\Complexesstack{\mathcal{C}\!\mathit{omplexes}}
% \Pic is the operator that assigns to X its picard group, usage \Pic(X)
% \Picardstack_{X/B} denotes the Picard stack of X over B
% \Picardfunctor_{X/B} denotes the Picard functor of X over B
\def\Pic{\mathop{\mathrm{Pic}}\nolimits}
\def\Picardstack{\mathcal{P}\!\mathit{ic}}
\def\Picardfunctor{\mathrm{Pic}}
\def\Deformationcategory{\mathcal{D}\!\mathit{ef}}

% --- appended 2026-10-12 (Chief): the project own definitions of the reference/slogan/history
% display environments, copied mechanically from the shared Stacks preamble lines 8-12
% (\usepackage{verbatim} plus the three \newenvironment lines, which the chapter preamble keeps
% verbatim from the project preamble). Nothing added, removed or redefined: these environments
% carry no mathematical content and the project itself typesets them as comments. Without these
% three lines a counted statement whose own text uses them fails with Environment slogan undefined.
\usepackage{verbatim}
\newenvironment{reference}{\comment}{\endcomment}
\newenvironment{slogan}{\comment}{\endcomment}
\newenvironment{history}{\comment}{\endcomment}
\usepackage{cleveref}
\crefname{originallemma}{original lemma}{original lemmas}
\Crefname{originallemma}{Original lemma}{Original lemmas}
\makeatletter
\newlabel{lemma-olsson-laszlo}{{0D6S}{source}{}{sourcepointer-1}{}}
\newlabel{lemma-olsson-laszlo@cref}{{[originallemma][0][]0D6S}{[1][0][]source}{}{}{}}
\makeatother
\crefname{originallemma}{original lemma}{original lemmas}
\Crefname{originallemma}{Original lemma}{Original lemmas}
\makeatletter
\newlabel{lemma-Rf-commutes-with-Rlim}{{0A07}{source}{}{sourcepointer-2}{}}
\newlabel{lemma-Rf-commutes-with-Rlim@cref}{{[originallemma][0][]0A07}{[1][0][]source}{}{}{}}
\makeatother
\crefname{originallemma}{original lemma}{original lemmas}
\Crefname{originallemma}{Original lemma}{Original lemmas}
\makeatletter
\newlabel{lemma-RGamma-commutes-with-Rlim}{{0D6K}{source}{}{sourcepointer-3}{}}
\newlabel{lemma-RGamma-commutes-with-Rlim@cref}{{[originallemma][0][]0D6K}{[1][0][]source}{}{}{}}
\makeatother
\crefname{originallemma}{original lemma}{original lemmas}
\Crefname{originallemma}{Original lemma}{Original lemmas}
\makeatletter
\newlabel{lemma-cohomology-derived-limit-injective}{{0D6L}{source}{}{sourcepointer-4}{}}
\newlabel{lemma-cohomology-derived-limit-injective@cref}{{[originallemma][0][]0D6L}{[1][0][]source}{}{}{}}
\makeatother
\crefname{originallemma}{original lemma}{original lemmas}
\Crefname{originallemma}{Original lemma}{Original lemmas}
\makeatletter
\newlabel{derived-lemma-two-ss-complex-functor}{{015J}{source}{}{sourcepointer-5}{}}
\newlabel{derived-lemma-two-ss-complex-functor@cref}{{[originallemma][0][]015J}{[1][0][]source}{}{}{}}
\makeatother
\crefname{originalremark}{original remark}{original remarks}
\Crefname{originalremark}{Original remark}{Original remarks}
\makeatletter
\newlabel{derived-remark-truncation-distinguished-triangle}{{08J5}{source}{}{sourcepointer-6}{}}
\newlabel{derived-remark-truncation-distinguished-triangle@cref}{{[originalremark][0][]08J5}{[1][0][]source}{}{}{}}
\makeatother
\title{A second version of the Olsson-Laszlo boundedness lemma for a morphism of ringed sites}
\author{The Stacks Project Authors}
\date{stacks-project:sites-cohomology.tex:local-snapshot-20261011; GNU FDL 1.2 or later}
\begin{document}
\maketitle
\noindent\textit{Editorial note.} All delivered mathematical text is the author's own, copied from the pinned Stacks Project chapter file sites-cohomology.tex; no statement or proof step was rewritten, reordered or removed.\par\medskip
\noindent\textit{Editorial note.} The proof is the author's complete argument: the derived limit presentation, the eventual constancy of the cohomology groups on a cofinal system of coverings, the identification of the cohomology of the derived limit and the final truncation comparison are all written out, with every cited input and the dimension bound named.\par\medskip
\noindent\textit{Editorial note.} Every labelled reference occurring inside the counted spans is carried in assembly.external\_references as a non-load-bearing navigation pointer (same-chapter ones at their definition site, cross-chapter ones at the invocation inside the counted span), so no load-bearing mathematics is replaced by a stub. No bibliography citation occurs inside the counted spans: the author's statement environment carries, inside a 'reference' environment at source lines 5877-5880, a provenance remark citing a published lemma of Olsson and Laszlo, and that remark is not part of the mathematical statement; it is therefore not delivered, and the counted statement is split into the two spans 5875-5876 and 5881-5896. No citation was supplied, invented or reconstructed, and no other text of the author was dropped.\par\medskip
\noindent\textit{Editorial note.} The checked chapter list for same-name collisions is the set of stacks chapters staged in this repository: chow, cohomology, divisors, etale-cohomology, flat, more-algebra, more-morphisms, morphisms, spaces-more-morphisms, varieties together with this card's chapter sites-cohomology.\par\medskip
\noindent\textit{Editorial note.} The preamble is the chapter's own retained minimal preamble (sources/stacks-sites-cohomology/retained-preamble.tex), a verbatim excerpt of the native preamble: the Xy-pic setup and the native theorem declarations with the mathematical macros.\par\medskip
\noindent\textit{Allowed standard prerequisites (editorial).}\begin{itemize}
\item The two Olsson-Laszlo situations. Same chapter, Situations situation-olsson-laszlo and situation-olsson-laszlo-prime. Supply the hypotheses referred to by assumptions (1) and (2).
\item Unbounded complexes are derived limits of their truncations. Same chapter, Lemma lemma-olsson-laszlo. Writes the complex as the derived limit of its canonical truncations.
\item The derived pushforward commutes with derived limits, and cohomology of a derived limit. Same chapter, Lemmas lemma-Rf-commutes-with-Rlim and lemma-RGamma-commutes-with-Rlim. Move the derived limit through the pushforward and identify the values of the resulting cohomology sheaves.
\item Constancy of a derived limit of a constant system. Same chapter, Lemma lemma-cohomology-derived-limit-injective. Turns the constancy of the system of cohomology groups into the asserted isomorphism.
\end{itemize}\medskip
\ifdefined\hypertarget\hypertarget{sourcepointer-1}{}\fi
\noindent\textit{Source reference.} Original lemma 0D6S: sources/stacks-sites-cohomology/pinned-original-source.tex, source lines 5707--5707. Navigation pointer invoked at proof line 5826; the counted span contains the pointer only and the cited result is not reproduced inside it, so the pointer is registered as a non-load-bearing navigation stub in the established form for Stacks material. No bibliography citation is involved.\par\medskip
\ifdefined\hypertarget\hypertarget{sourcepointer-2}{}\fi
\noindent\textit{Source reference.} Original lemma 0A07: sources/stacks-sites-cohomology/pinned-original-source.tex, source lines 5150--5150. Navigation pointer invoked at proof line 5827; the counted span contains the pointer only and the cited result is not reproduced inside it, so the pointer is registered as a non-load-bearing navigation stub in the established form for Stacks material. No bibliography citation is involved.\par\medskip
\ifdefined\hypertarget\hypertarget{sourcepointer-3}{}\fi
\noindent\textit{Source reference.} Original lemma 0D6K: sources/stacks-sites-cohomology/pinned-original-source.tex, source lines 5126--5126. Navigation pointer invoked at proof line 5791; the counted span contains the pointer only and the cited result is not reproduced inside it, so the pointer is registered as a non-load-bearing navigation stub in the established form for Stacks material. No bibliography citation is involved.\par\medskip
\ifdefined\hypertarget\hypertarget{sourcepointer-4}{}\fi
\noindent\textit{Source reference.} Original lemma 0D6L: sources/stacks-sites-cohomology/pinned-original-source.tex, source lines 5279--5279. Navigation pointer invoked at proof line 5759; the counted span contains the pointer only and the cited result is not reproduced inside it, so the pointer is registered as a non-load-bearing navigation stub in the established form for Stacks material. No bibliography citation is involved.\par\medskip
\ifdefined\hypertarget\hypertarget{sourcepointer-5}{}\fi
\noindent\textit{Source reference.} Original lemma 015J: sources/stacks-sites-cohomology/pinned-original-source.tex, source lines 5742--5742. Navigation pointer invoked at proof line 5742; the counted span contains the pointer only and the cited result is not reproduced inside it, so the pointer is registered as a non-load-bearing navigation stub in the established form for Stacks material. No bibliography citation is involved.\par\medskip
\ifdefined\hypertarget\hypertarget{sourcepointer-6}{}\fi
\noindent\textit{Source reference.} Original remark 08J5: sources/stacks-sites-cohomology/pinned-original-source.tex, source lines 5948--5948. Navigation pointer invoked at proof line 5948; the counted span contains the pointer only and the cited result is not reproduced inside it, so the pointer is registered as a non-load-bearing navigation stub in the established form for Stacks material. No bibliography citation is involved.\par\medskip
\begin{situation}
\label{situation-olsson-laszlo}
Let $\mathcal{C}$ be a site. Let $\mathcal{O}$ be a sheaf of rings on
$\mathcal{C}$. Let $\mathcal{A} \subset \textit{Mod}(\mathcal{O})$
be a weak Serre subcategory. We assume the following is true:
there exists a subset $\mathcal{B} \subset \Ob(\mathcal{C})$ such that
\begin{enumerate}
\item every object of $\mathcal{C}$ has a covering whose
members are in $\mathcal{B}$, and
\item for every $V \in \mathcal{B}$ there exists an integer $d_V$
and a cofinal system $\text{Cov}_V$ of coverings of $V$ such
that
$$
H^p(V_i, \mathcal{F}) = 0 \text{ for }
\{V_i \to V\} \in \text{Cov}_V,\ p > d_V, \text{ and }
\mathcal{F} \in \Ob(\mathcal{A})
$$
\end{enumerate}
\end{situation}

\begin{situation}
\label{situation-olsson-laszlo-prime}
Let
$f : (\mathcal{C}, \mathcal{O}) \to (\mathcal{C}', \mathcal{O}')$
be a morphism of ringed sites. Let $u : \mathcal{C}' \to \mathcal{C}$
be the corresponding continuous functor of sites.
Let $\mathcal{A} \subset \textit{Mod}(\mathcal{O})$
be a weak Serre subcategory. We assume the following is true:
there exists a subset $\mathcal{B}' \subset \Ob(\mathcal{C}')$ such that
\begin{enumerate}
\item every object of $\mathcal{C}'$ has a covering whose
members are in $\mathcal{B}'$, and
\item for every $V' \in \mathcal{B}'$ there exists an integer $d_{V'}$
and a cofinal system $\text{Cov}_{V'}$ of coverings of $V'$ such
that
$$
H^p(u(V'_i), \mathcal{F}) = 0 \text{ for }
\{V'_i \to V'\} \in \text{Cov}_{V'},\ p > d_{V'}, \text{ and }
\mathcal{F} \in \Ob(\mathcal{A})
$$
\end{enumerate}
\end{situation}

\begin{lemma}
\label{lemma-olsson-laszlo-map-version-two}

Let $f : (\mathcal{C}, \mathcal{O}) \to (\mathcal{C}', \mathcal{O}')$
be a morphism of ringed sites.
assume moreover there is an integer $N$ such that
\begin{enumerate}
\item $\mathcal{C}, \mathcal{O}, \mathcal{A}$ satisfy the
assumption of Situation \ref{situation-olsson-laszlo},
\item $f : (\mathcal{C}, \mathcal{O}) \to (\mathcal{C}', \mathcal{O}')$
and $\mathcal{A}$ satisfy the assumption of
Situation \ref{situation-olsson-laszlo-prime},
\item $R^pf_*\mathcal{F} = 0$ for
$p > N$ and $\mathcal{F} \in \Ob(\mathcal{A})$,
\end{enumerate}
Then for $K$ in $D_\mathcal{A}(\mathcal{O})$ the map
$H^j(Rf_*K) \to H^j(Rf_*(\tau_{\geq -n}K))$ is an isomorphism
for $j \geq N - n$.
\end{lemma}

\begin{proof}
Let $K$ be in $D_\mathcal{A}(\mathcal{O})$.
By Lemma \ref{lemma-olsson-laszlo} we have $K = R\lim \tau_{\geq -n}K$.
By Lemma \ref{lemma-Rf-commutes-with-Rlim}
we have $Rf_*K = R\lim Rf_*(\tau_{\geq -n}K)$.
Let $V' \in \mathcal{B}'$ and let $\{V'_i \to V'\}$ be an element
of $\text{Cov}_{V'}$. Then we consider
$$
H^j(V'_i, Rf_*K) = H^j(u(V'_i), K)
\quad\text{and}\quad
H^j(V'_i, Rf_*(\tau_{\geq -n}K)) = H^j(u(V'_i), \tau_{\geq -n}K)
$$
The assumption in Situation \ref{situation-olsson-laszlo-prime}
implies that the last group is independent of $n$ for $n$ large enough
depending on $j$ and $d_{V'}$. Some details omitted.
We apply this for $j$ and $j - 1$ and via
Lemma \ref{lemma-RGamma-commutes-with-Rlim} this gives that
$$
H^j(V'_i, Rf_*K) = \lim H^j(V'_i, Rf_*(\tau_{\geq -n} K))
$$
and the system on the right is constant for $n$ larger than
a constant depending only on $d_{V'}$ and $j$.
Thus Lemma \ref{lemma-cohomology-derived-limit-injective}
implies that
$$
H^j(Rf_*K)(V') \longrightarrow
\left(\lim H^j(Rf_*(\tau_{\geq -n}K))\right)(V')
$$
is injective. Since the elements $V' \in \mathcal{B}'$ cover
every object of $\mathcal{C}'$ we conclude that the map
$H^j(Rf_*K) \to \lim H^j(Rf_*(\tau_{\geq -n}K))$ is injective.
The spectral sequence
$$
E_2^{p, q} = R^pf_*H^q(\tau_{\geq -n}K)
$$
converging to $H^{p + q}(Rf_*(\tau_{\geq -n}K))$
(Derived Categories, Lemma \ref{derived-lemma-two-ss-complex-functor})
and assumption (3) show that
$H^j(Rf_*(\tau_{\geq -n}K))$ is constant for $n \geq N - j$.
Hence $H^j(Rf_*K) \to H^j(Rf_*(\tau_{\geq -n}K))$ is injective
for $j \geq N - n$.

\medskip\noindent
Thus we proved the lemma with ``isomorphism'' in the last line of
the lemma replaced by ``injective''. However, now choose $j$ and $n$
with $j \geq N - n$. Then consider the distinguished triangle
$$
\tau_{\leq -n - 1}K \to K \to \tau_{\geq -n}K \to (\tau_{\leq -n - 1}K)[1]
$$
See Derived Categories, Remark
\ref{derived-remark-truncation-distinguished-triangle}.
Since $\tau_{\geq -n}\tau_{\leq -n -1}K = 0$, the
injectivity already proven for $\tau_{-n - 1}K$ implies
$$
0 = H^j(Rf_*(\tau_{\leq -n - 1}K)) =
H^{j + 1}(Rf_*(\tau_{\leq -n - 1}K)) =
H^{j + 2}(Rf_*(\tau_{\leq -n - 1}K)) = \ldots
$$
By the long exact cohomology sequence associated to the distinguished
triangle
$$
Rf_*(\tau_{\leq -n - 1}K) \to Rf_*K \to Rf_*(\tau_{\geq -n}K) \to
Rf_*(\tau_{\leq -n - 1}K)[1]
$$
this implies that $H^j(Rf_*K) \to H^j(Rf_*(\tau_{\geq -n}K))$
is an isomorphism.
\end{proof}

\end{document}
